\documentclass[11pt]{article}
\usepackage{amsmath,amssymb,geometry,booktabs,xcolor}
\geometry{margin=2.5cm}
\title{The Koide relation from two Casimir ratios:\\
  $Q = C_2(\wedge^2\!5)/C_2(5)$, $\delta = C_2(3)/C_2(S^3\!3)$}
\author{A.~Rivero}
\date{March 2026 -- working draft v11}

\begin{document}
\maketitle

\begin{abstract}
Both the Koide ratio $Q=3/2$ and the Koide angle
$\delta=2/9$ are identified as quadratic Casimir ratios
of simple Lie groups.
(1)~$Q = C_2(\wedge^2\mathbf{5})/C_2(\mathbf{5}) = 3/2$
in SU(5)$_{\rm flavour}$ (the sBootstrap turtle group),
uniquely selecting $N_f=5$.
(2)~$\delta = C_2(\mathbf{3})/C_2(S^3\mathbf{3}) = (4/3)/6 = 2/9$
in SU(3)$_{\rm generation}$ (three generations embedded in
$G_2 = {\rm Aut}(\mathbb{O})$).
The chain angle factor $\delta_{scb} \approx 3\delta_\ell$
is the colour multiplicity $N_c=3$.
Together, these determine the charged lepton mass ratios
from group theory alone: $Q$ constrains the trace ratio
(how much masses spread) and $\delta$ constrains the angular
position (where they sit on the Koide orbit).
The sBootstrap provides the physical context: SU(5)$_{\rm flavour}$
acts on the five turtle quarks, SU(3)$_{\rm generation}$ acts
on the three generations, and $N_c=3$ is QCD colour.
No new parameters are introduced.
\end{abstract}

\section{The two Casimir identities}

\subsection{$Q = 3/2$ from SU(5)}

The Koide ratio $Q \equiv (\sum\sqrt{m_i})^2/\sum m_i$ equals
the quadratic Casimir ratio of the antisymmetric two-tensor
to the fundamental of SU($N$):
\begin{equation}
  Q = \frac{C_2(\wedge^2\mathbf{N})}{C_2(\mathbf{N})}
  = \frac{2(N-2)}{N-1}.
\end{equation}
Setting $Q = 3/2$: $N = 5$.  In the sBootstrap, SU(5) is the
flavour group of the five turtle quarks $(u,c,d,s,b)$.

Further SU(5) invariants give $Z$-ladder integers:
$T(\mathbf{15})/T(\mathbf{5}) = N+2 = 7$,\;
$T(\mathbf{24})/T(\mathbf{5}) = 2N = 10$.

\subsection{$\delta = 2/9$ from SU(3) $\subset$ $G_2$}

The Koide angle $\delta = 2/9$ equals the Casimir ratio
of the fundamental to the symmetric three-tensor of SU(3):
\begin{equation}
  \delta = \frac{C_2(\mathbf{3}_{SU(3)})}{C_2(S^3\mathbf{3}_{SU(3)})}
  = \frac{4/3}{6} = \frac{2}{9}.
\end{equation}
Here SU(3) is the \emph{generation} symmetry (three generations
of fermions), naturally embedded in $G_2 = {\rm Aut}(\mathbb{O})$
(the automorphism group of the octonions, which provides the
exceptional Jordan algebra $J_3(\mathbb{O})$ as a three-generation
mass matrix framework).

The symmetric three-tensor $S^3(\mathbf{3})$ has
dimension $\binom{5}{3} = 10$ and Casimir
$C_2 = p(N+p)(N-1)/(2N)|_{N=3,p=3} = 6$.

\section{The two groups}

\begin{center}
\begin{tabular}{@{}l l l l@{}}
\toprule
Group & Acts on & Casimir ratio & Result \\
\midrule
SU(5)$_{\rm flavour}$ & 5 turtle quarks
  & $C_2(\wedge^2)/C_2(\rm fund) = 3/2$ & $Q_{\rm Koide}$ \\
SU(3)$_{\rm generation}$ & 3 generations
  & $C_2({\rm fund})/C_2(S^3) = 2/9$ & $\delta_{\rm Koide}$ \\
\bottomrule
\end{tabular}
\end{center}

\noindent
The two groups act on \emph{different spaces}: SU(5)
mixes flavours ($u,c,d,s,b$), SU(3) mixes generations
($1,2,3$).  Their Casimir ratios determine orthogonal
aspects of the mass spectrum: $Q$ (the trace condition)
and $\delta$ (the angular position).

\section{The chain angle as colour multiplicity}

The chain angle $\delta_{scb} \approx 37.2^\circ \approx 0.649$
is close to $3\delta_\ell = 3 \times 2/9 = 2/3 \approx 38.2^\circ$
(deviation $1.0^\circ$).  The factor $3$ is $N_c$, the number of
QCD colours:
\begin{equation}
  \delta_{scb} \approx N_c \cdot \delta_\ell.
\end{equation}
This connects the quark Koide chain to the lepton Koide through
the colour multiplicity, completing the group-theoretic picture:
\begin{equation}
  Q \leftarrow \text{SU}(5)_{\rm flav},\qquad
  \delta \leftarrow \text{SU}(3)_{\rm gen},\qquad
  \delta_{scb}/\delta_\ell \leftarrow \text{SU}(3)_{\rm colour}.
\end{equation}

\section{What is and is not derived}

\begin{center}\small
\begin{tabular}{@{}l l l@{}}
\toprule
Quantity & Formula & Status \\
\midrule
$Q = 3/2$ & $C_2(\wedge^2\!\mathbf{5})/C_2(\mathbf{5})$
  & Group theory (exact) \\
$\delta = 2/9$ & $C_2(\mathbf{3})/C_2(S^3\!\mathbf{3})$
  & Group theory (exact) \\
$N = 5$ & $2(N-2)/(N-1) = 3/2$ & Uniquely selected \\
$\delta_{scb}/\delta_\ell = 3$ & $N_c$
  & Colour multiplicity ($1^\circ$ off) \\
\midrule
\emph{Why} mass matrix $\sim$ Casimirs & --- & Not derived \\
$Z$-ladder formula & --- & Not derived \\
Overall mass scale $M_0$ & --- & Free parameter \\
\bottomrule
\end{tabular}
\end{center}

\noindent
The Casimir identities provide the \emph{values} of $Q$ and
$\delta$ from group theory, and \emph{select} $N=5$ and
$N_{\rm gen}=3$ as the only viable integers.  The
\emph{mechanism} that forces the mass matrix to satisfy these
Casimir ratios remains the key open question.  Candidates
include: a conformal fixed point where anomalous dimensions are
Casimir-scaled; a vacuum alignment condition in the Higgs
potential; or a string-theoretic construction where the
octonion structure of $G_2$ determines the mass matrix.

\begin{thebibliography}{5}
\bibitem{Koide1983}
Y.~Koide, Phys.\ Rev.\ D \textbf{28}, 252 (1983).
\bibitem{Rivero2024}
A.~Rivero, Eur.\ Phys.\ J.\ C \textbf{84}, 1058 (2024).
\end{thebibliography}

\end{document}
